\section{Estudio de mercado y viabilidad}\label{sec:03-mercado}

\subsection{Mercado y tendencias}

El turismo internacional cerró 2025 con 1\,520 millones de llegadas (+4\,\%); Europa recibió
793 millones \cite{mk:untourism2025}.
España batió su récord con 96,8 millones de turistas internacionales y 134\,712~M€ de gasto
\cite{mk:egatur2025}. Las OTA facturaron 94\,000~M\$ en 2024 (+7\,\% anual) y ven
en la IA conversacional una amenaza \cite{mk:skiftota}. Los ingresos se generan al \emph{reservar};
\emph{planificar}, el eslabón que cubre Kyrian World, apenas se monetiza por sí mismo.

En España el 40\,\% de los adultos ya ha usado un asistente de IA para preparar un viaje (el
73\,\% entre los 18 y los 24 años) \cite{mk:yougov2026}; Europa va por detrás (Reino Unido
22\,\%, Francia 19\,\%, Alemania 15\,\%) \cite{mk:phocuseurope2025}. Falta confianza: solo el 17--20\,\% de los viajeros europeos confía plenamente en una IA
\cite{mk:phocustrust}, y el 49\,\% de los españoles que no la usan teme recomendaciones
inexactas o desactualizadas \cite{mk:yougov2026}. En una encuesta patrocinada por Mindtrip a
103 organizaciones de gestión de destinos (DMO) de EE.\,UU., solo el 22\,\% ofrece un
planificador con IA en su web \cite{mk:dmoreadiness}.

\subsection{Competencia}

La \cref{tab:03-competidores} enseña que los generadores sin datos propios acaban comprados y
que Expedia retiró Romie por responder sin inventario ni precios reales \cite{mk:romie}.

\begin{table*}[tb]\centering\footnotesize
\begin{tabularx}{\textwidth}{@{}p{1.9cm}p{2.4cm}YYY@{}}\toprule
Nombre & Tipo / propietario & Negocio & Diferenciador & Debilidad o estado \\\midrule
Layla & \emph{Startup} ($\approx$5~M€) & \emph{Freemium} (9,99~\$/mes), afiliación \cite{mk:laylapricing} & Chat y vídeo; precios en vivo & Comprada por Expedia (07-2026); mantiene su marca \cite{mk:laylaexpedia} \\
Roam Around & \emph{Startup} (1~M\$) & Gratuito & Sencillez; auge tras ChatGPT & Absorbida por Layla (02-2024); sin datos propios \\
Mindtrip & \emph{Startup} (19~M\$) & Gratis; comisión; B2B & 11~M de POI; >35 destinos clientes \cite{mk:mindtrip} & Fuerte en EE.\,UU.; poco español \\
GuideGeek & Matador Network & Marca blanca para DMO & >70 DMO; WhatsApp, Instagram \cite{mk:guidegeek} & Depende de OpenAI; generalista \\
Wonderplan & \emph{Startup} & Gratis, sin registro & Rapidez; PDF & Casi insensible al perfil \\
Wanderlog & \emph{Startup} (YC) & \emph{Freemium} & Colaboración y mapas & IA secundaria \\
ChatGPT & OpenAI & Suscripción; apps de Booking y Expedia (10-2025) \cite{mk:chatgptapps} & Distribución masiva & Sin anclaje; reserva agéntica experimental \\
Google AI Mode & Alphabet & Publicidad, comisión & Itinerarios (11-2025); datos de Maps \cite{mk:googleaimode} & Solo EE.\,UU.; Google Trips cerró en 2019 \\
Expedia Romie & Expedia Group & Comisión OTA & Planificación y soporte (2024) & Retirado: sin inventario ni precios reales \cite{mk:romie} \\
Booking.com & Booking Holdings & Comisión OTA & Chat sobre inventario propio (2023) & Ligado a su inventario; app en ChatGPT \\\bottomrule
\end{tabularx}
\caption{Competidores y comparables (septiembre de 2026). Las cifras sin cita proceden del
estudio de mercado del proyecto (reseñas y prensa sectorial).}
\label{tab:03-competidores}
\end{table*}

\textbf{Posicionamiento.} Kyrian World compite en \emph{verificabilidad}, no en distribución ni
inventario: cada tarjeta remite, con enlace a la fuente, a un documento de un corpus con licencia
limpia (\cref{sec:05-arquitectura-ia}). Es un anclaje documental (existencia, descripción, fuente),
no de disponibilidad ni precio: cubrir la carencia de Romie exige integrar inventario (horizonte
de 12 meses). Asume una ciudad por viaje (\adr{0019}) y pocas ciudades, aunque en castellano. Su
ventaja es el coste fijo: $\approx$5~€/mes en recursos de Terraform (\adr{0023}) más
$\approx$8~\$/mes de un WAF creado a mano, unos 12--13~€/mes (estimación; la factura de septiembre
debe sustituirla), frente a rondas de millones.

\subsection{Monetización}

La \textbf{afiliación} paga por reserva completada (\cref{tab:03-afiliacion}), pero el modelo por sesión de Booking.com solo computa la reserva hecha en el
navegador inicial \cite{mk:bookingcj}. La \textbf{suscripción} toma como referencia Layla
Premium, a 9,99~\$/mes o 49,99~\$/año \cite{mk:laylapricing}; se supone un plan de 39~€/año con
una conversión del 1--2\,\% (\emph{hipótesis}: el contenido del plan \emph{premium} está sin definir). La vía con más recorrido es el \textbf{B2B o
marca blanca para DMO}, con demanda demostrada por GuideGeek \cite{mk:guidegeek} y Mindtrip
\cite{mk:mindtrip}: se vendería el \emph{pipeline} \code{/add-city}, a 6\,000--15\,000~€ por
destino y año (\emph{hipótesis}).

\begin{table}[tb]\centering\footnotesize
\begin{tabularx}{\columnwidth}{@{}lY@{}}\toprule
Programa & Comisión de referencia \\\midrule
Booking.com & 4\,\% de estancias completadas; por sesión \cite{mk:bookingcj} \\
Expedia Group & Hoteles 3--4\,\%; actividades 4\,\% \cite{mk:expediaaff} \\
GetYourGuide / Viator & $\approx$8\,\%; \emph{cookie} de 30 días \cite{mk:viator} \\
Kiwi.com & 3\,\% del billete \cite{mk:kiwi} \\
Omio & 2--8\,\% según mercado y volumen \cite{mk:omio} \\\bottomrule
\end{tabularx}
\caption{Comisiones de afiliación de referencia (consultadas el 23-09-2026).}
\label{tab:03-afiliacion}
\end{table}

\subsection{Economía unitaria}

Con unas 12 llamadas al modelo por viaje, $\approx$50\,000 \emph{tokens} de entrada y
$\approx$8\,000 de salida (\emph{estimación}), el coste variable por viaje es de 0,09~\$ con
Claude Haiku 4.5, el modelo de producción; 0,18--0,30~\$ con Sonnet (el \emph{endpoint} de la UE
es un 10\,\% más caro) \cite{mk:anthropicpricing}; 0,07~\$ con Nova Pro \cite{mk:bedrockpricing};
vectores y \emph{embeddings}, $<$0,001~\$. El escenario usa 0,20~€/viaje para dejar margen a
reintentos, llamadas de chat y un posible modelo mayor.

El ingreso por afiliación es de 0,37~€ por viaje en el escenario base (0,15--0,76~€), con tasas
de clic y conversión \emph{supuestas}: el margen queda en $\approx$0,17~€ por viaje, y en el
escenario pesimista cada viaje pierde dinero. En B2C, el LTV ronda 1,8~€ (\emph{hipótesis}:
1,8 viajes/año $\times$ 0,17~€ más 1,5\,\% $\times$ 39~€ de \emph{premium}, $\approx$0,9~€/año
durante dos años). Como CAC se usa de \emph{proxy} el coste por instalación de una app móvil de
viajes, 4,20~\$ \cite{mk:cac}: LTV/CAC $\approx$0,5, así que el B2C depende de canales orgánicos. En B2B, un
contrato de 8\,000~€/año durante tres años frente a 2\,000--4\,000~€ de coste de venta pública da
un LTV/CAC de 6--12 (\emph{hipótesis}).

\begin{table}[tb]\centering\footnotesize
\begin{tabularx}{\columnwidth}{@{}Yrrr@{}}\toprule
Concepto (€) & Año 1 & Año 2 & Año 3 \\\midrule
Viajes planificados & 5\,200 & 40\,000 & 162\,000 \\
Afiliación (0,37~€/viaje) & 1\,924 & 14\,800 & 59\,940 \\
\emph{Premium} (1--2\,\% $\times$ 39~€) & 1\,560 & 14\,625 & 70\,200 \\
B2B (contratos) & 3\,000 (1) & 32\,000 (4) & 108\,000 (12) \\
\textbf{Ingresos} & \textbf{6\,484} & \textbf{61\,425} & \textbf{238\,140} \\
\quad peso del B2B & 46\,\% & 52\,\% & 45\,\% \\\midrule
Coste variable (0,20~€/viaje) & 1\,040 & 8\,000 & 32\,400 \\
Infra., datos y legal & 1\,980 & 7\,440 & 19\,800 \\
Marketing y venta B2B & 3\,000 & 28\,000 & 95\,000 \\
\textbf{Resultado antes de salarios} & \textbf{+464} & \textbf{+17\,985} & \textbf{+90\,940} \\
Salarios & 0 & 35\,000 & 90\,000 \\
\textbf{Resultado neto} & \textbf{+464} & \textbf{$-$17\,015} & \textbf{+940} \\\bottomrule
\end{tabularx}
\caption{Escenario base a tres años (10-2026 a 09-2029). Hipótesis: 4\,000 / 25\,000 / 90\,000
usuarios activos con 1,3 / 1,6 / 1,8 viajes al año; B2B a 3\,000~€ (piloto), 8\,000~€ y
9\,000~€ por contrato; fundador asalariado desde el año 2 y una segunda persona en el año 3.}
\label{tab:03-escenario}
\end{table}

En la \cref{tab:03-escenario} el año~1 cuadra solo porque no imputa salario: con un salario
mínimo imputado daría $\approx-$15\,000~€ (\emph{hipótesis}). Cubrir 35\,000~€ de salario y 5\,000~€ de fijos
exigiría, solo con afiliación, unos 235\,000 viajes al año, irreal sin adquisición de pago; con
B2B bastan cinco contratos de $\approx$8\,000~€.

\subsection{Marco legal}

Entre corchetes, el estado de cada obligación. El \textbf{RGPD} exige base jurídica,
minimización, plazos de conservación de las trazas y cuidado con el texto libre, que puede revelar
categorías especiales \cite{mk:anthropicpricing}; [\emph{pendiente}: el login menciona una política de privacidad y unos
términos cuyas páginas no existen]. El \textbf{Reglamento de IA} (UE) 2024/1689 no incluye un
planificador de viajes en el Anexo~III: le afecta el art.~50 (desde el 02-08-2026): avisar de que
se habla con una IA y marcar de forma legible por máquina el contenido sintético \cite{mk:aiact50} [\emph{parcial}: el planner se
titula ``Dile a la IA adónde quieres ir'', sin aviso explícito ni marcado]. La \textbf{Directiva de
viajes combinados} no aplica mientras Kyrian World solo redirija. Las \textbf{licencias} imponen
atribución: CC BY-SA 4.0 (Wikipedia, Wikivoyage) obliga a compartir igual las adaptaciones
publicadas, como el corpus JSONL \cite{ai:wikivoyagelicense}; ODbL exige atribuir a OpenStreetMap \cite{ai:osmcopyright}
[\emph{parcial}: la tarjeta enlaza la fuente, pero la licencia solo aparece en un \emph{tooltip};
la introducción de la ciudad sí cita ``Wikivoyage (CC BY-SA 4.0)'']. La API gratuita de
\textbf{Open-Meteo} es solo no comercial \cite{mk:openmeteo} [\emph{pendiente}: monetizar exige su
plan comercial u otra fuente].

\textbf{Vender a una DMO pública española.} La licencia hipotética de 6\,000--15\,000~€ encaja en
un contrato menor de servicios (art.~118 LCSP, $<$15\,000~€) \cite{fx:lcsp}. Exigiría además
conformidad con el Esquema Nacional de Seguridad \cite{fx:ens}, sin evaluar; accesibilidad según el
RD~1112/2018 y la EN~301~549 \cite{fx:rd1112,fx:en301549}, declarada pero no auditada
(\cref{sec:02-producto}); alojamiento en la UE, que la región \code{eu-west-1} cubre para la API
pero sin acreditación formal; contenido en español, cuando el 98,8\,\% del corpus está en inglés, y
cobertura de ciudades medianas, donde Bolonia tiene solo 1\,597 documentos. Hoy ninguno está
acreditado.

\begin{table}[tb]\centering\footnotesize
\begin{tabularx}{\columnwidth}{@{}YY@{}}\toprule
\textbf{Fortalezas} & \textbf{Debilidades} \\\midrule
\begin{itemize}[leftmargin=*,label=--,nosep]
\item Anclaje citado, licencia limpia
\item 12--13~€/mes de coste fijo
\item Alta de ciudad reproducible
\end{itemize} &
\begin{itemize}[leftmargin=*,label=--,nosep]
\item Pocas ciudades; sin reservas
\item Sin métricas ni ingresos
\item \emph{Bus factor} 1
\end{itemize} \\\midrule
\textbf{Oportunidades} & \textbf{Amenazas} \\\midrule
\begin{itemize}[leftmargin=*,label=--,nosep]
\item 78\,\% de DMO sin IA (EE.\,UU.)
\item Red DTI: 440 destinos
\item Art.~50; LLM más barato
\end{itemize} &
\begin{itemize}[leftmargin=*,label=--,nosep]
\item Google y ChatGPT gratis
\item Expedia compra Layla
\item Comisiones y datos abiertos
\end{itemize} \\\bottomrule
\end{tabularx}
\caption{DAFO de Kyrian World (septiembre de 2026).}
\label{tab:03-dafo}
\end{table}

\subsection{Conclusión de viabilidad}

Como B2C financiado con afiliación (\cref{tab:03-dafo}), la viabilidad es marginal: el ingreso
por viaje es del orden del coste del modelo, el LTV no cubre el CAC y las plataformas regalan la planificación. El nicho recomendado
es el B2B para oficinas de turismo de ciudades europeas medianas y los 440 destinos de la Red de
Destinos Turísticos Inteligentes \cite{mk:segittur}, con el viajero hispanohablante como segmento
secundario. El diferenciador a defender es la \emph{IA verificable con fuentes}, en línea con la
transparencia del art.~50; la carencia de Romie \cite{mk:romie}, inventario y precios reales, sigue
sin cubrir.
